\scanpage{67}
\mspage{60}
% archive sheet number: 66
b) Tout ens. ordonné fini est admissible. Plus généralement (sauf erreur) tt ens. ordonné $I$ tel que l'espace topologique associé \unsure{à} $I^{\unsure{\circ|o}}$ soit \unsure{noeth.} (i.e. \deleted{toute} suite croiss. \unsure{d'ouverts} de $I^\circ$ est stationnaire, ou encore tte suite croissante de fermés de $I$ est stationnaire), et si de plus tte partie fermée irréductible de $I^\circ$ admet un pt générique, i.e. $I^\circ$ est (non seulement un espace noethérien, mais aussi) un espace top. \uline{sobre}. (Pour $I$, cela signifie que \deleted{\unsure{toute}} partie ouverte filtrante croissante de $I$ admet un plus grand élément\unsure{,|.}\deleted{\unsure{)}})

\deleted{\unsure{Ceci implique}}
c) En particulier, si $I',I$ \deleted{\unsure{font des} $\mathrm{Diag}$.} sont des ens. ordonnés finis, \unsure{toute flèche|les hom} $u:I'\unsure{-|\to}I$ est admissible.
\inserted{$=$, plus généralement, si $I^\circ,I'^\circ$ sont des \deleted{\unsure{espaces}} \inserted{ens. ordonnés} \deleted{\unsure{noeth. et sobres}} et $I'^\circ$ est un espace noeth. sobre.}

2) Supposons $u:I'\unsure{-|\to}I$ admissible, \deleted{\unsure{alors}} \deleted{\unsure{précisons que les} $\unsure{u_*|I}$} considérons
\diagram{
$u_{\unsure{*|\circ}}^W:$
\begin{tikzcd}[ampersand replacement=\&,column sep=huge,row sep=large]
M(I')_fW(I')_f^{-1} \arrow[r,dash] \arrow[d,"\unsure{q}"',"\approx"] \& \deleted{\unsure{M}}M(I)_fW(I)_f^{-1}\arrow[d,"\approx"] \\
\mathrm{HO}_W(I') \& \mathrm{HO}_W(I)
\end{tikzcd}
}
induit par $u_*$. Ce foncteur est adjoint à droite de
\[
u_W^*:\mathrm{HO}_W(I)\unsure{-|\longrightarrow}\mathrm{HO}_W(I').
\]
